\documentclass[a4paper]{article}
\usepackage[fontset=fandol]{ctex}
\usepackage[margin=1.65cm]{geometry}
\usepackage{amsmath,amssymb}
\usepackage{graphicx}
\usepackage{hyperref}
\usepackage{longtable,booktabs}
\usepackage{xcolor}
\usepackage{enumitem}
\usepackage{listings}
\lstset{basicstyle=\ttfamily\small,breaklines=true,columns=fullflexible,keepspaces=true}
\hypersetup{colorlinks=true,linkcolor=blue,urlcolor=blue}
\setlist[itemize]{leftmargin=1.25em,itemsep=0.18em}
\setlength{\parskip}{0.45em}
\setlength{\parindent}{2em}
\sloppy
\title{第 07 讲：多 GPU 并行基础}
\author{CS336 Spring 2026 public videos and official course materials}
\date{}
\begin{document}
\maketitle
\tableofcontents
\newpage
\section{本章导读}
本章讨论 \textbf{多 GPU 并行基础}。CS336 的第一层目标不是把现成大模型当作黑箱使用，而是把 language model 从输入文本、训练目标、模型结构、数据、系统和评估这些部件重新拆开。只有拆开之后，读者才能判断一个设计选择究竟改变了什么。
设想单张 GPU 已经放不下模型，或者单卡训练时间不可接受。并行训练不是简单把 GPU 数乘上去，而是要决定哪些张量复制、哪些切分、哪些通过 all-reduce、all-gather 或 reduce-scatter 通信。
本章对应 CS336 Spring 2026 第 07 个公开视频，并结合课程页提供的可解析资料。视频和课程页给出事实边界；正文把这些材料改写为可自学的教材章节。
阅读本章时，先理解问题为什么存在，再看定义、公式和实现形态，随后通过例题和练习检查自己是否能独立复现这条 reasoning chain。
\begin{figure}[h]
\centering
\includegraphics[width=0.58\linewidth]{source_anchor.jpg}
\caption{第 07 讲的课程图像锚点。}
\end{figure}
\subsection{本章知识路线}
\begin{itemize}
\item 本章首先说明 多 GPU 并行基础 为什么是 language modeling pipeline 中不可跳过的一环。
\item 其次定义本章反复使用的术语，并说明这些术语对应的计算对象或系统对象。
\item 随后用公式和伪代码表达主要机制，让概念能够落到可计算的变量上。
\item 最后给出例题、风险讨论和练习，便于读者检查自己是否真正掌握了机制。
\end{itemize}
\subsection{结构图}
下面的图用于帮助读者定位本章的核心结构。先看概念之间的依赖，再回到正文中的公式、伪代码和例题。
\begin{figure}[h]
\centering
\includegraphics[width=0.72\linewidth]{figures/source_grounded_concept_map.png}
\caption{第 07 讲概念图：本章问题、术语和机制的关系。}
\end{figure}
\begin{figure}[h]
\centering
\includegraphics[width=0.72\linewidth]{figures/source_grounded_checklist.png}
\caption{第 07 讲学习路径：从问题定义进入公式、实现、实验和练习。}
\end{figure}
\section{本章主线}
多 GPU 并行基础 的主线是把 language model 的抽象计算放回真实机器：哪些张量要移动，哪些状态要保存，哪些瓶颈会在训练或服务时显现。
\subsection{本章要解决的核心问题}
本讲把上一周的单 GPU 层级扩展为多 GPU/多节点层级：NVLink、NVSwitch、InfiniBand、Ethernet 和 RDMA 都会进入通信账本。
读完这一段后，应能用一两句话说明 多 GPU 并行基础 要解决的瓶颈，并指出这个瓶颈发生在数据、模型、优化、系统还是评估层。
\subsection{从机制到资源账本}
collective communication 是并行训练的语言。broadcast、scatter、gather、reduce、all-gather、reduce-scatter、all-reduce、all-to-all 分别对应不同的数据移动模式。
随后把机制写成账本：输入规模是什么，主要状态是什么，成本或质量指标怎样随变量改变。这个账本让后面的公式不再只是符号。
\subsection{从课堂讲解到可复现实验}
data parallel、tensor parallel、pipeline parallel 和 sequence/expert parallel 的选择取决于模型能否放下、互连是否足够快、通信能否和计算重叠。
最后把课堂讲解转成可复现实验：固定协议，改变一个变量，记录中间量，并解释结果为何支持或反驳原来的机制判断。
\section{术语、符号与问题边界}
本章的术语横跨模型、数据和系统。读者需要先知道每个词指向的对象：有些词描述数学目标，有些词描述张量形状，有些词描述硬件瓶颈，还有一些词描述评估协议。术语边界不清，后面的公式和实现就会变成空洞的缩写。
\begin{longtable}{p{0.08\linewidth}p{0.42\linewidth}p{0.42\linewidth}}
\toprule
\textbf{\#} & \textbf{术语} & \textbf{阅读要求}\\
\midrule
1 & data parallelism（数据并行） & 说明它如何改变训练样本或 token 分布，并给出一个会改变结论的数据边界条件。\\
2 & tensor parallelism（张量并行） & 给出定义、一个最小例子、一个关联公式或代码位置，以及一个典型失败情形。\\
3 & pipeline parallelism（流水线并行） & 给出定义、一个最小例子、一个关联公式或代码位置，以及一个典型失败情形。\\
4 & all-reduce & 给出定义、一个最小例子、一个关联公式或代码位置，以及一个典型失败情形。\\
5 & communication cost（通信成本） & 给出定义、一个最小例子、一个关联公式或代码位置，以及一个典型失败情形。\\
\bottomrule
\end{longtable}
这些术语之间有依赖关系。例如 tokenization 改变 sequence length，sequence length 又改变 attention FLOPs、KV cache 和 evaluation token budget；parallelism 改变显存占用，也改变通信瓶颈和故障恢复方式。
\subsection{术语之间的关系}
在本章中，data parallelism（数据并行） 给出问题入口，tensor parallelism（张量并行） 描述主要机制，pipeline parallelism（流水线并行） 则把机制连接到可测量的结果。读教材时应当沿着这个链条追问：输入是什么，中间状态怎样改变，输出指标为什么随之变化。
课程材料还会反复出现 data、rank、gpus、parallelism、e.g.、ranks 等词。它们不是一组同义标签，而是提示本章的不同层次：有的指对象，有的指操作，有的指约束或指标。
因此，术语表不是背诵清单，而是后文阅读的索引。遇到一个术语时，先定位它属于 representation、optimization、systems、data 还是 evaluation，再判断它改变哪一项账本。
\section{为什么单 GPU 不够}
本节讨论 为什么单 GPU 不够。它先限定 多 GPU 并行基础 要回答的基本问题：哪些现象需要解释，哪些变量可以测量，哪些结论必须回到 吞吐、延迟、显存峰值、通信量和 kernel 利用率 中检验。
\subsection{课程目标与问题边界}
视频从上一周单 GPU kernel 优化过渡到多 GPU：目标是利用多张 GPU 训练更大模型或更快完成训练。
并行不是免费加速；参数、activation、optimizer state、gradient 和数据都可能需要复制、切分或通信。
选择并行策略要同时满足显存容量、计算吞吐和互连带宽约束。
系统部分的核心问题是把数学计算映射到硬件。一个算法的复杂度公式只能给出方向；真正的速度还取决于内存层级、通信拓扑、kernel shape 和调度方式。
\subsection{为什么单 GPU 不够 的机制}
视频从上一周单 GPU kernel 优化过渡到多 GPU：目标是利用多张 GPU 训练更大模型或更快完成训练。 在 为什么单 GPU 不够 中，核心对象是硬件和系统状态：同一算法会因为 memory hierarchy、通信拓扑或 kernel shape 不同而表现出完全不同的速度。
并行不是免费加速；参数、activation、optimizer state、gradient 和数据都可能需要复制、切分或通信。 因而同一个方法在不同规模下可能改变不同的变量，尤其是吞吐、延迟、显存峰值、通信量和 kernel 利用率。比较方法时，必须同时给出问题规模和实验条件。
选择并行策略要同时满足显存容量、计算吞吐和互连带宽约束。 可检验性来自 profile trace、roofline 估算和端到端延迟分解；如果这些证据互相矛盾，就要先解释 measurement setup。
因此，为什么单 GPU 不够 的学习目标是建立一张成本地图：读者应能指出瓶颈在计算、内存、通信还是调度，并能设计一个小 benchmark 去验证。
\subsection{资源、效率与效果的关系}
\[\text{speedup}(n)=\frac{T_1}{T_n},\qquad \text{efficiency}(n)=\frac{\text{speedup}(n)}{n}\]
符号说明：\(T_1\) 是单设备时间，\(T_n\) 是 \(n\) 设备时间；efficiency 衡量扩展损失。
这个关系式应被读成 为什么单 GPU 不够 的资源账本。它不承诺精确预测每次运行时间，而是帮助判断主要成本来自计算、内存访问、通信还是调度。
在系统实验中，tokens/s、TTFT、显存峰值和 kernel time 都是 proxy。它们各自只观察一部分系统，因此报告时要同时写清 workload shape、warmup、同步方式和硬件型号。
\subsection{把课程目标写成检查流程}
把 为什么单 GPU 不够 写成程序时，最重要的是暴露状态变化和数据移动。下面的伪代码保留执行顺序，省略框架样板。
\begin{lstlisting}[language=Python]
if model_fits_each_gpu:
    start_with_data_parallel()
else:
    shard_parameters_and_activations()
\end{lstlisting}
阅读这段实现时，重点看 为什么单 GPU 不够 怎样移动张量、占用显存并触发通信。实验记录至少要包含 吞吐、延迟、显存峰值、通信量和 kernel 利用率，否则很难判断瓶颈来自算子、内存层级还是调度。
\subsection{边界条件与常见误解}
\begin{itemize}
\item 能跑起来不等于高效；通信可能吞掉全部扩展收益。
\item 并行策略常和 batch size、sequence length、optimizer state 绑定。
\item 小规模 kernel benchmark 不一定预测端到端训练或 serving latency，因为调度、通信和请求形状会改变瓶颈。
\item 硬件规格表和框架 profiler 都需要结合 workload 解释，不能把单个峰值数字当作机制证明。
\end{itemize}
\subsection{与本章其他概念的关系}
为什么单 GPU 不够 连接本章的数学机制和工程成本：术语表说明对象，公式估算资源，伪代码暴露状态变化，例题则让读者检查一次完整的成本推理。
\subsection{怎样检验 这一课程目标}
检验 这一课程目标 时，先把抽象说法落到硬件和系统状态上。与 parallelism、multi GPU、speedup、efficiency、communication 相关的对象必须能被构造、打印、计数或评分；否则实验只能得到描述性观察，不能得到可复现结论。
一个合适的小实验应当保留本节机制，同时让 吞吐、延迟、显存峰值、通信量和 kernel 利用率 中至少一个量发生可解释变化。输入规模不必逼真；关键是读者能手算或打印关键中间量，并说明哪个变量触发了结果变化。
随后把公式说明转成可记录的数量关系。符号说明：\textbackslash{}(T\_1\textbackslash{}) 是单设备时间，\textbackslash{}(T\_n\textbackslash{}) 是 \textbackslash{}(n\textbackslash{}) 设备时间；efficiency 衡量扩展损失。 如果数量只能间接观测，就必须说明使用了什么 proxy；如果使用 benchmark 或 reward，也要说明协议和随机性。
系统消融通常一次只改变 batch shape、sequence length、GPU 数、kernel 实现或调度策略。若同时改变硬件和 workload，实验只能说明端到端表现不同，不能定位瓶颈。
最后检查失败条件。工程判断往往在反例中变清楚：能跑起来不等于高效；通信可能吞掉全部扩展收益；并行策略常和 batch size、sequence length、optimizer state 绑定。复习时可以把这些限制改写成反例测试，观察它们怎样改变 吞吐、延迟、显存峰值、通信量和 kernel 利用率。
把实验结论放回 多 GPU 并行基础 时，要说明 为什么单 GPU 不够 改变的是训练吞吐、推理延迟、显存峰值还是通信开销。能做出这种归因，才说明读者已经从接口使用进入系统解释。
\subsection{概念辨析：为什么单 GPU 不够}
parallelism、multi GPU、speedup、efficiency、communication 常一起出现在系统讲解中，但它们分别指硬件对象、程序操作、瓶颈来源或测量指标。读者应先定位每个词在执行路径中的位置。
其中，parallelism（并行）把模型、数据、序列或专家切分到多设备；GPU 是 CS336 systems 部分的主要硬件对象。
区分这些词时，可以先问它们是否改变 吞吐、延迟、显存峰值、通信量和 kernel 利用率。一个术语如果不能对应到变量、状态、代码路径或评估协议，就还没有形成可操作含义。
围绕 为什么单 GPU 不够 复习时，可以写一个小 benchmark：固定输入形状，改变一个系统变量，记录 profile trace 和端到端时间。
\subsection{为什么单 GPU 不够 的小结}
\begin{itemize}
\item 讨论 为什么单 GPU 不够 时，应同时报告问题规模、实验设置和主要观测变量，尤其是 吞吐、延迟、显存峰值、通信量和 kernel 利用率。
\item 如果方法声称更快，要区分计算减少、内存访问减少、通信隐藏和调度改善；这些机制不能用一个 tokens/s 数字互相替代。
\item profile trace、roofline 估算和端到端 benchmark 应互相校验，单独一个指标不足以解释系统瓶颈。
\end{itemize}
\section{Data parallelism 与 all-reduce}
现在进入 Data parallelism 与 all-reduce。这一节把算法放到硬件和运行时中考察：同一个数学表达式在不同 batch、shape、GPU 拓扑和 kernel 实现下会得到不同的成本账本。
\subsection{为什么要讨论 Data parallelism 与 all-reduce}
data parallelism 复制模型，把 batch 切给不同 GPU，反向传播后聚合梯度。
all-reduce 是核心通信原语：每个设备最终得到全局梯度和。
梯度累积可以在显存有限时模拟更大 batch，也可以减少通信频率，但会改变优化动态。
数据部分的核心问题是分布控制。采集、过滤、去重和混合并不会只改变文件大小，它们会改变模型看到的语言、知识、偏差和 benchmark contamination 风险。
\subsection{Data parallelism 与 all-reduce 的机制}
data parallelism 复制模型，把 batch 切给不同 GPU，反向传播后聚合梯度。 在 Data parallelism 与 all-reduce 中，核心对象是硬件和系统状态：同一算法会因为 memory hierarchy、通信拓扑或 kernel shape 不同而表现出完全不同的速度。
all-reduce 是核心通信原语：每个设备最终得到全局梯度和。 因而同一个方法在不同规模下可能改变不同的变量，尤其是吞吐、延迟、显存峰值、通信量和 kernel 利用率。比较方法时，必须同时给出问题规模和实验条件。
梯度累积可以在显存有限时模拟更大 batch，也可以减少通信频率，但会改变优化动态。 可检验性来自 profile trace、roofline 估算和端到端延迟分解；如果这些证据互相矛盾，就要先解释 measurement setup。
因此，Data parallelism 与 all-reduce 的学习目标是建立一张成本地图：读者应能指出瓶颈在计算、内存、通信还是调度，并能设计一个小 benchmark 去验证。
\subsection{Data parallelism 与 all-reduce 的资源账本}
\[g=\frac{1}{n}\sum_{i=1}^{n}g_i,\qquad \theta\leftarrow \theta-\eta g\]
符号说明：\(g_i\) 是第 \(i\) 个 worker 的局部梯度；all-reduce 后每个 worker 使用同一个全局梯度更新。
这个关系式应被读成 Data parallelism 与 all-reduce 的资源账本。它不承诺精确预测每次运行时间，而是帮助判断主要成本来自计算、内存访问、通信还是调度。
在系统实验中，tokens/s、TTFT、显存峰值和 kernel time 都是 proxy。它们各自只观察一部分系统，因此报告时要同时写清 workload shape、warmup、同步方式和硬件型号。
\subsection{Data parallelism 与 all-reduce 在数据管线中的实现}
把 Data parallelism 与 all-reduce 写成程序时，最重要的是暴露状态变化和数据移动。下面的伪代码保留执行顺序，省略框架样板。
\begin{lstlisting}[language=Python]
loss.backward()
all_reduce(param.grad, op='sum')
param.grad /= world_size
optimizer.step()
\end{lstlisting}
阅读这段实现时，重点看 Data parallelism 与 all-reduce 怎样移动张量、占用显存并触发通信。实验记录至少要包含 吞吐、延迟、显存峰值、通信量和 kernel 利用率，否则很难判断瓶颈来自算子、内存层级还是调度。
\subsection{边界条件与常见误解}
\begin{itemize}
\item global batch 变大后 learning rate schedule 可能需要调整。
\item all-reduce 成本与参数量和网络拓扑有关，不只与 GPU 数量有关。
\item 小规模 kernel benchmark 不一定预测端到端训练或 serving latency，因为调度、通信和请求形状会改变瓶颈。
\item 硬件规格表和框架 profiler 都需要结合 workload 解释，不能把单个峰值数字当作机制证明。
\end{itemize}
\subsection{与本章其他概念的关系}
Data parallelism 与 all-reduce 连接本章的数学机制和工程成本：术语表说明对象，公式估算资源，伪代码暴露状态变化，例题则让读者检查一次完整的成本推理。
\subsection{怎样检验 Data parallelism 与 all-reduce}
检验 Data parallelism 与 all-reduce 时，先把抽象说法落到硬件和系统状态上。与 data parallel、all-reduce、gradient、batch、optimizer 相关的对象必须能被构造、打印、计数或评分；否则实验只能得到描述性观察，不能得到可复现结论。
一个合适的小实验应当保留本节机制，同时让 吞吐、延迟、显存峰值、通信量和 kernel 利用率 中至少一个量发生可解释变化。输入规模不必逼真；关键是读者能手算或打印关键中间量，并说明哪个变量触发了结果变化。
随后把公式说明转成可记录的数量关系。符号说明：\textbackslash{}(g\_i\textbackslash{}) 是第 \textbackslash{}(i\textbackslash{}) 个 worker 的局部梯度；all-reduce 后每个 worker 使用同一个全局梯度更新。 如果数量只能间接观测，就必须说明使用了什么 proxy；如果使用 benchmark 或 reward，也要说明协议和随机性。
系统消融通常一次只改变 batch shape、sequence length、GPU 数、kernel 实现或调度策略。若同时改变硬件和 workload，实验只能说明端到端表现不同，不能定位瓶颈。
最后检查失败条件。工程判断往往在反例中变清楚：global batch 变大后 learning rate schedule 可能需要调整；all-reduce 成本与参数量和网络拓扑有关，不只与 GPU 数量有关。复习时可以把这些限制改写成反例测试，观察它们怎样改变 吞吐、延迟、显存峰值、通信量和 kernel 利用率。
把实验结论放回 多 GPU 并行基础 时，要说明 Data parallelism 与 all-reduce 改变的是训练吞吐、推理延迟、显存峰值还是通信开销。能做出这种归因，才说明读者已经从接口使用进入系统解释。
\subsection{概念辨析：Data parallelism 与 all-reduce}
data parallel、all-reduce、gradient、batch、optimizer 常一起出现在系统讲解中，但它们分别指硬件对象、程序操作、瓶颈来源或测量指标。读者应先定位每个词在执行路径中的位置。
其中，data（数据）是模型能力、偏差和法律风险的主要来源；all-reduce 聚合各 rank 梯度，是 data parallel training 的核心 collective。
区分这些词时，可以先问它们是否改变 吞吐、延迟、显存峰值、通信量和 kernel 利用率。一个术语如果不能对应到变量、状态、代码路径或评估协议，就还没有形成可操作含义。
围绕 Data parallelism 与 all-reduce 复习时，可以写一个小 benchmark：固定输入形状，改变一个系统变量，记录 profile trace 和端到端时间。
\subsection{Data parallelism 与 all-reduce 的小结}
\begin{itemize}
\item 讨论 Data parallelism 与 all-reduce 时，应同时报告问题规模、实验设置和主要观测变量，尤其是 吞吐、延迟、显存峰值、通信量和 kernel 利用率。
\item 如果方法声称更快，要区分计算减少、内存访问减少、通信隐藏和调度改善；这些机制不能用一个 tokens/s 数字互相替代。
\item profile trace、roofline 估算和端到端 benchmark 应互相校验，单独一个指标不足以解释系统瓶颈。
\end{itemize}
\section{Tensor/pipeline parallelism 的切分直觉}
现在进入 Tensor/pipeline parallelism 的切分直觉。这一节把算法放到硬件和运行时中考察：同一个数学表达式在不同 batch、shape、GPU 拓扑和 kernel 实现下会得到不同的成本账本。
\subsection{为什么要讨论 Tensor/pipeline parallelism 的切分直觉}
tensor parallelism 切矩阵或 head，让单层计算分布到多设备；pipeline parallelism 切层，把不同 micro-batches 送入流水线。
pipeline bubble 是流水线并行的典型低效来源；micro-batching 是主要缓解手段。
真正的训练系统常混合 data、tensor、pipeline、sequence/context parallelism。
系统部分的核心问题是把数学计算映射到硬件。一个算法的复杂度公式只能给出方向；真正的速度还取决于内存层级、通信拓扑、kernel shape 和调度方式。
\subsection{Tensor/pipeline parallelism 的切分直觉 的机制}
tensor parallelism 切矩阵或 head，让单层计算分布到多设备；pipeline parallelism 切层，把不同 micro-batches 送入流水线。 在 Tensor/pipeline parallelism 的切分直觉 中，核心对象是硬件和系统状态：同一算法会因为 memory hierarchy、通信拓扑或 kernel shape 不同而表现出完全不同的速度。
pipeline bubble 是流水线并行的典型低效来源；micro-batching 是主要缓解手段。 因而同一个方法在不同规模下可能改变不同的变量，尤其是吞吐、延迟、显存峰值、通信量和 kernel 利用率。比较方法时，必须同时给出问题规模和实验条件。
真正的训练系统常混合 data、tensor、pipeline、sequence/context parallelism。 可检验性来自 profile trace、roofline 估算和端到端延迟分解；如果这些证据互相矛盾，就要先解释 measurement setup。
因此，Tensor/pipeline parallelism 的切分直觉 的学习目标是建立一张成本地图：读者应能指出瓶颈在计算、内存、通信还是调度，并能设计一个小 benchmark 去验证。
\subsection{Tensor/pipeline parallelism 的切分直觉 的资源账本}
\[\text{pipeline bubble fraction}\approx\frac{p-1}{m+p-1}\]
符号说明：\(p\) 是 pipeline stages，\(m\) 是 micro-batches；micro-batches 越多，bubble 比例越低。
这个关系式应被读成 Tensor/pipeline parallelism 的切分直觉 的资源账本。它不承诺精确预测每次运行时间，而是帮助判断主要成本来自计算、内存访问、通信还是调度。
在系统实验中，tokens/s、TTFT、显存峰值和 kernel time 都是 proxy。它们各自只观察一部分系统，因此报告时要同时写清 workload shape、warmup、同步方式和硬件型号。
\subsection{从 Tensor/pipeline parallelism 的切分直觉 到程序}
把 Tensor/pipeline parallelism 的切分直觉 写成程序时，最重要的是暴露状态变化和数据移动。下面的伪代码保留执行顺序，省略框架样板。
\begin{lstlisting}[language=Python]
split_layers_across_stages(model)
for microbatch in schedule:
    forward_stage_then_backward_stage(microbatch)
\end{lstlisting}
阅读这段实现时，重点看 Tensor/pipeline parallelism 的切分直觉 怎样移动张量、占用显存并触发通信。实验记录至少要包含 吞吐、延迟、显存峰值、通信量和 kernel 利用率，否则很难判断瓶颈来自算子、内存层级还是调度。
\subsection{边界条件与常见误解}
\begin{itemize}
\item micro-batch 增多会增加 activation 管理复杂度。
\item tensor parallelism 需要高带宽互连，否则 matmul 中间结果通信会成为瓶颈。
\item 小规模 kernel benchmark 不一定预测端到端训练或 serving latency，因为调度、通信和请求形状会改变瓶颈。
\item 硬件规格表和框架 profiler 都需要结合 workload 解释，不能把单个峰值数字当作机制证明。
\end{itemize}
\subsection{与本章其他概念的关系}
Tensor/pipeline parallelism 的切分直觉 连接本章的数学机制和工程成本：术语表说明对象，公式估算资源，伪代码暴露状态变化，例题则让读者检查一次完整的成本推理。
\subsection{怎样检验 Tensor/pipeline parallelism 的切分直觉}
检验 Tensor/pipeline parallelism 的切分直觉 时，先把抽象说法落到硬件和系统状态上。与 tensor parallel、pipeline、microbatch、bubble、stage 相关的对象必须能被构造、打印、计数或评分；否则实验只能得到描述性观察，不能得到可复现结论。
一个合适的小实验应当保留本节机制，同时让 吞吐、延迟、显存峰值、通信量和 kernel 利用率 中至少一个量发生可解释变化。输入规模不必逼真；关键是读者能手算或打印关键中间量，并说明哪个变量触发了结果变化。
随后把公式说明转成可记录的数量关系。符号说明：\textbackslash{}(p\textbackslash{}) 是 pipeline stages，\textbackslash{}(m\textbackslash{}) 是 micro-batches；micro-batches 越多，bubble 比例越低。 如果数量只能间接观测，就必须说明使用了什么 proxy；如果使用 benchmark 或 reward，也要说明协议和随机性。
系统消融通常一次只改变 batch shape、sequence length、GPU 数、kernel 实现或调度策略。若同时改变硬件和 workload，实验只能说明端到端表现不同，不能定位瓶颈。
最后检查失败条件。工程判断往往在反例中变清楚：micro-batch 增多会增加 activation 管理复杂度；tensor parallelism 需要高带宽互连，否则 matmul 中间结果通信会成为瓶颈。复习时可以把这些限制改写成反例测试，观察它们怎样改变 吞吐、延迟、显存峰值、通信量和 kernel 利用率。
把实验结论放回 多 GPU 并行基础 时，要说明 Tensor/pipeline parallelism 的切分直觉 改变的是训练吞吐、推理延迟、显存峰值还是通信开销。能做出这种归因，才说明读者已经从接口使用进入系统解释。
\subsection{概念辨析：Tensor/pipeline parallelism 的切分直觉}
tensor parallel、pipeline、microbatch、bubble、stage 常一起出现在系统讲解中，但它们分别指硬件对象、程序操作、瓶颈来源或测量指标。读者应先定位每个词在执行路径中的位置。
其中，tensor（张量）是参数、激活、梯度和数据的共同表示。
区分这些词时，可以先问它们是否改变 吞吐、延迟、显存峰值、通信量和 kernel 利用率。一个术语如果不能对应到变量、状态、代码路径或评估协议，就还没有形成可操作含义。
围绕 Tensor/pipeline parallelism 的切分直觉 复习时，可以写一个小 benchmark：固定输入形状，改变一个系统变量，记录 profile trace 和端到端时间。
\subsection{Tensor/pipeline parallelism 的切分直觉 的小结}
\begin{itemize}
\item 讨论 Tensor/pipeline parallelism 的切分直觉 时，应同时报告问题规模、实验设置和主要观测变量，尤其是 吞吐、延迟、显存峰值、通信量和 kernel 利用率。
\item 如果方法声称更快，要区分计算减少、内存访问减少、通信隐藏和调度改善；这些机制不能用一个 tokens/s 数字互相替代。
\item profile trace、roofline 估算和端到端 benchmark 应互相校验，单独一个指标不足以解释系统瓶颈。
\end{itemize}
\section{综合讨论：从局部机制到完整系统}
本章的主题是 多 GPU 并行基础，但它不应被读成几个相邻知识点的拼接。为什么单 GPU 不够、Data parallelism 与 all-reduce、Tensor/pipeline parallelism 的切分直觉 共同回答同一个问题：语言模型系统中哪些对象可以被定义，哪些成本可以被估算，哪些结论必须通过实验或评估来确认。
这些对象包括 data parallelism（数据并行）、tensor parallelism（张量并行）、pipeline parallelism（流水线并行）、all-reduce、communication cost（通信成本）。读者可以把它们看成一组接口：有的接口连接文本和 token，有的连接模型结构和张量计算，有的连接数据样本和训练目标，有的连接离线评估和在线服务。接口一旦改变，后面的成本、错误类型和优化空间也会随之改变。
\subsection{为什么单 GPU 不够 在系统中的位置}
为什么单 GPU 不够 首先提供一个局部解释：视频从上一周单 GPU kernel 优化过渡到多 GPU：目标是利用多张 GPU 训练更大模型或更快完成训练。 这类解释的价值在于把直觉压缩成可操作对象。一旦对象明确，读者就可以追问它的输入、输出、中间状态和失败条件，而不是只记住方法名。
其次，为什么单 GPU 不够 会改变至少一个资源或评估变量。它可能改变 sequence length、activation size、memory bandwidth、communication volume、data mixture、reward distribution 或 benchmark protocol。这些变量在小例子里可以手算，在真实训练和服务系统里则需要 profiling、logging 和 held-out evaluation。
最后，为什么单 GPU 不够 不能脱离边界条件理解。一个关键 caveat 是：能跑起来不等于高效；通信可能吞掉全部扩展收益。。把 caveat 写出来不是为了削弱结论，而是为了说明结论在哪些输入、硬件、数据和评估条件下才成立。
\subsection{Data parallelism 与 all-reduce 在系统中的位置}
Data parallelism 与 all-reduce 首先提供一个局部解释：data parallelism 复制模型，把 batch 切给不同 GPU，反向传播后聚合梯度。 这类解释的价值在于把直觉压缩成可操作对象。一旦对象明确，读者就可以追问它的输入、输出、中间状态和失败条件，而不是只记住方法名。
其次，Data parallelism 与 all-reduce 会改变至少一个资源或评估变量。它可能改变 sequence length、activation size、memory bandwidth、communication volume、data mixture、reward distribution 或 benchmark protocol。这些变量在小例子里可以手算，在真实训练和服务系统里则需要 profiling、logging 和 held-out evaluation。
最后，Data parallelism 与 all-reduce 不能脱离边界条件理解。一个关键 caveat 是：global batch 变大后 learning rate schedule 可能需要调整。。把 caveat 写出来不是为了削弱结论，而是为了说明结论在哪些输入、硬件、数据和评估条件下才成立。
\subsection{Tensor/pipeline parallelism 的切分直觉 在系统中的位置}
Tensor/pipeline parallelism 的切分直觉 首先提供一个局部解释：tensor parallelism 切矩阵或 head，让单层计算分布到多设备；pipeline parallelism 切层，把不同 micro-batches 送入流水线。 这类解释的价值在于把直觉压缩成可操作对象。一旦对象明确，读者就可以追问它的输入、输出、中间状态和失败条件，而不是只记住方法名。
其次，Tensor/pipeline parallelism 的切分直觉 会改变至少一个资源或评估变量。它可能改变 sequence length、activation size、memory bandwidth、communication volume、data mixture、reward distribution 或 benchmark protocol。这些变量在小例子里可以手算，在真实训练和服务系统里则需要 profiling、logging 和 held-out evaluation。
最后，Tensor/pipeline parallelism 的切分直觉 不能脱离边界条件理解。一个关键 caveat 是：micro-batch 增多会增加 activation 管理复杂度。。把 caveat 写出来不是为了削弱结论，而是为了说明结论在哪些输入、硬件、数据和评估条件下才成立。
\subsection{把本章知识用于读论文和复现实验}
读相关论文时，可以把本章变成一个检查表。第一，论文是否清楚定义了输入规模、模型规模、数据来源和评估协议；第二，论文声称的改进落在哪个变量上；第三，作者是否报告了会让结论失效的设置；第四，实验是否能分离模型结构、数据质量、优化设置和系统实现的影响。
复现实验时，也应避免直接追求大规模。先用最小程序复现一个公式、一个伪代码分支或一个 profile 现象，再扩大输入规模。这样做的原因很简单：如果小规模程序无法解释中间量，大规模训练只会把错误隐藏在日志和随机波动里。
把这些检查合在一起，本章给出的不是一个固定答案，而是一种阅读 CS336 的方法：每个模型机制都要落到变量，每个变量都要落到测量，每个测量都要说明协议，每个协议都要承认边界。
因此，本章中的任何一句结论都可以被改写成一个可引用的完整句式：在给定数据、模型、硬件和评估协议下，某个机制改变了某个变量，并通过某个观测指标表现出来。这样的句子比单独背诵术语更长，但它包含了自学、复现和引用时真正需要的信息。
把本章用于自学时，可以按三遍阅读。第一遍只追踪对象和定义，确认每个术语到底指向文本、token、张量、样本、请求还是评估记录。第二遍追踪公式和伪代码，确认每一步改变了哪些状态。第三遍追踪实验和 caveat，确认哪些结论只在特定规模、数据或硬件条件下成立。
把本章用于复习时，则应反过来操作：先从一个失败案例或异常指标出发，再回到相应机制。loss 曲线异常可能指向数据、优化或模型容量；吞吐异常可能指向 kernel、通信或调度；benchmark 异常可能指向 contamination、prompt protocol 或采样设置。能够从异常反推机制，是掌握本章的实用标准。
\section{例题：把概念变成可计算检查}
\subsection{Data parallel all-reduce 的通信账本}
每张 GPU 得到局部梯度 `g\_i`，训练需要全局平均梯度。
\begin{itemize}
\item all-reduce 让每个 rank 都得到 `sum\_i g\_i`。
\item 再除以 world size，得到全局 batch 的平均梯度。
\item 参数越大、通信网络越慢，扩展效率越容易下降。
\end{itemize}
\textbf{结论。} 数据并行的数学等价很简单，系统代价集中在梯度通信和 batch/LR 调整。
这个例题的目的不是替代完整工程实现，而是给出一种最小推理形式：把问题转成变量，写出近似公式或伪代码，再说明近似何时失效。
\section{实验设计：从小例子到可复现结论}
学习 多 GPU 并行基础 不能停在概念层。一个教材化结论至少需要四件事：可构造的输入、可观测的中间量、可比较的指标，以及清楚的失败条件。下面的实验设计不是要求训练大模型，而是要求读者用小规模程序复现本章机制。
\subsection{最小输入}
先构造只触发本章核心机制的小输入。例如 tokenizer 章节可以用几个包含 rare words 和 Unicode 字符的字符串；GPU 章节可以用一个 elementwise kernel 和一个 GEMM；data 章节可以用十几篇近重复文档。小输入的价值在于它让错误可见。
\subsection{资源账本}
对同一输入写下 FLOPs、bytes moved、显存峰值、通信量或样本数量的估算。估算不需要一开始就精确，但必须说明单位和近似假设。随后用 profiler、benchmark、validation loss 或人工检查去验证估算是否同量级。
\subsection{单变量消融}
一次只改变一个因素，例如 vocabulary size、head 数、batch size、filter threshold、reward scale、decoding temperature 或 GPU 数量。若多个变量同时变化，实验只能说明系统变了，不能说明机制是什么。
\subsection{失败条件}
最后要刻意触发失败：极端 context length、错误 dtype、过小 batch、污染 benchmark、低质量数据或不可靠 verifier。失败条件常常比成功曲线更能说明一个方法的边界。
\subsection{记录与报告}
实验记录应能让另一个读者复现同一结论。最少要写清楚输入构造、代码版本、随机种子、硬件、batch shape、数据版本、评价指标和异常处理方式。报告结论时不要只写“更快”或“更好”，而要说明快在哪里、好在哪个指标上，以及是否牺牲了内存、稳定性、数据质量或可解释性。
常见报告错误是把局部现象写成普遍规律。例如一次小模型实验里的 loss 改善，不等于更大模型或不同数据分布上也会改善；一次 kernel benchmark 的吞吐提升，不等于端到端 serving latency 一定降低。教材中的实验报告应始终把结论限定在可复现条件内。
\subsection{从小实验到大结论}
小实验的作用不是替代课程中的完整训练或系统实现，而是建立一条可检查的推理链。先在小输入上确认变量方向，再在中等规模上确认数量级，最后才讨论是否能外推到更大模型、更多 token 或更复杂 workload。缺少中间层级时，大结论通常只是在复述直觉。
因此，实验报告最好同时包含正例和反例：正例说明机制在受控条件下怎样工作，反例说明条件稍变时哪里失效。对 CS336 这样的课程，反例尤其重要，因为 language model 的错误常来自多个层面叠加：数据分布、优化稳定性、硬件瓶颈、评估协议和服务负载都可能同时改变观察结果。
还要记录资料版本与复现边界。公开视频、课程页、配套代码和 PDF 可能在不同时间更新；复现实验应写明使用的材料版本，并说明哪些结论来自课程事实，哪些是为了帮助理解而加入的解释性推导。这样才能把学习笔记变成可检查的技术记录。
若复现实验与课程结论不一致，首先检查输入、版本和评估协议，而不是立即修改模型假设。
\subsection{把本章重新读成一条证据链}
完成小实验之后，可以把 多 GPU 并行基础 重新读成一条证据链：先用“为什么单 GPU 不够”确定问题入口，再用“Data parallelism 与 all-reduce”解释主要机制，最后用“Tensor/pipeline parallelism 的切分直觉”检查边界或外推条件。这种读法比按页背诵更接近教材训练，因为它要求每个结论都能说明前提、变量和证据。
本章反复使用的术语包括 data parallelism（数据并行）、tensor parallelism（张量并行）、pipeline parallelism（流水线并行）、all-reduce、communication cost（通信成本）。复习时不要只写定义，而要为每个术语补上一句操作性说明：它在课程代码、公式、数据记录或评估协议中对应什么对象；如果这个对象被删除、替换或缩放，哪一个观测量会变化。
最后，用一个反例结束复习。反例可以是资源估算不准、数据污染、reward 失真、benchmark 解析失败、长上下文显存爆炸或 scaling 外推失效。能解释反例的章节，才真正具备自学和引用价值；不能解释反例的章节，只是把课程材料换了一种排版。
引用本章结论时，还应保留来源边界：哪些说法直接来自视频、课程页、slides 或代码，哪些是为了串联教材叙述而加入的解释性推导。这个边界不会削弱结论，反而让读者知道何处可以复核，何处需要在自己的实验中重新验证。如果后续课程材料更新，也可以沿着同一证据链替换来源，而不必重写整章结构。
\section{课程资料与回看路径}
教材正文已经把主要概念、公式和实现逻辑组织成连续叙述；本节只提供回到课程材料的阅读路径。读者复习时可以先完成前文练习，再按下面的时间段或资料组核对细节。
\subsection{视频回看路径}
\begin{longtable}{p{0.20\linewidth}p{0.46\linewidth}p{0.26\linewidth}}
\toprule
\textbf{时间} & \textbf{关键词} & \textbf{复习任务}\\
\midrule
00:00:05-00:05:42 & going, single, we're, gpus, your, talk, then & 回看定义、例子、推导或 caveat 在该段中的位置。\\
00:19:02-00:25:39 & gpus, then, inaudible, rank, connected, particular, very & 回看定义、例子、推导或 caveat 在该段中的位置。\\
00:38:59-00:46:07 & going, then, reduce, it's, processes, rank, input & 回看定义、例子、推导或 caveat 在该段中的位置。\\
00:53:43-01:00:16 & going, data, then, it's, cuda, data., batch & 回看定义、例子、推导或 caveat 在该段中的位置。\\
01:14:14-01:20:44 & parallelism, then, data, tensor, pipeline, parallel, we'll & 回看定义、例子、推导或 caveat 在该段中的位置。\\
\bottomrule
\end{longtable}
这些时间段用于定位视频中的讲解顺序。严肃引用时，应回到视频片段确认讲者表述、板书或幻灯片内容。
回看视频时，不必逐字抄录字幕。更有效的方法是把每个时间段判定为定义、例子、推导、代码解释或 caveat，再把它放回本章对应小节。
\subsection{课程资料阅读路径}
\begin{longtable}{p{0.14\linewidth}p{0.48\linewidth}p{0.30\linewidth}}
\toprule
\textbf{资料组} & \textbf{标题/页块} & \textbf{关键词}\\
\midrule
official\_01 & Lecture 7: parallelism & single, week, parallelism, node, multi-gpu, last, across\\
official\_03 & - Data parallelism: DDP (all-reduce), FSDP/ZeRO (all-gather + reduce-scatter) & parallelism, interconnects, pipeline, work, reduce, store, memory\\
official\_05 & Use case: after backward pass, sum gradients from different data shards, but ... & data, rank, splits, case, backward, pass, gradients\\
official\_08 & Notes: & across, ranks, data, same, sharding, strategy, distributed\\
official\_10 & Function: pipeline\_parallelism & function, pipeline\_parallelism, pipeline\_parallelism\_main, setup, cleanup, spawn, get\_init\_params\\
\bottomrule
\end{longtable}
课程资料通常给出更稳定的标题、公式、图或代码结构；视频则补充动机和口头解释。两者合起来构成本章的事实边界。
阅读资料时，可以把每个资料组拆成三个对象：一个定义，一个公式或伪代码动作，一个边界条件。这样比逐字翻译页面或脚本更容易形成可用知识。
\section{实践说明、历史位置与风险}
多 GPU 并行基础 在 CS336 的课程结构中连接基础机制和规模化问题。基础机制解释方法为什么成立；规模化问题检验它在更大模型、更长上下文、更复杂数据或真实 serving workload 下是否仍成立。
历史位置不等于模型年表。更重要的是识别问题如何迁移：早期方法解决小规模建模问题，现代方法常常解决同一问题在大规模数据、GPU 集群、长上下文或人类偏好约束下的新形态。
\begin{itemize}
\item 资料版本限制（version risk）：课程页、公开视频和配套资料可能在学期中更新；引用时应注明访问日期或资料版本。
\item 测量口径限制（measurement risk）：FLOPs、tokens/s、benchmark score、reward 等数字只有在协议完整时可比较。
\item 规模外推限制（scaling risk）：小模型或小 batch 的机制不必然外推到 frontier-scale。
\item 数据分布限制（data risk）：数据来源、过滤和去重策略会改变模型行为，也会改变 evaluation contamination 风险。
\item 系统风险（systems risk）：性能剖析结果依赖硬件、driver、kernel、shape 和 batch distribution；脱离 workload 的性能数字不可直接复用。
\end{itemize}
\section{复习要点}
下面的问题把本章内容压缩为可反复检查的条目。每个条目都应能回到前文的解释、公式、伪代码或例题。
\begin{longtable}{p{0.16\linewidth}p{0.38\linewidth}p{0.38\linewidth}}
\toprule
\textbf{条目} & \textbf{对象} & \textbf{检查动作}\\
\midrule
条目 1 & data parallelism（数据并行） & 定义它；写出一个最小例子；说明一个 failure mode；指出它影响哪个资源或评估账本。\\
条目 2 & tensor parallelism（张量并行） & 定义它；写出一个最小例子；说明一个 failure mode；指出它影响哪个资源或评估账本。\\
条目 3 & pipeline parallelism（流水线并行） & 定义它；写出一个最小例子；说明一个 failure mode；指出它影响哪个资源或评估账本。\\
条目 4 & all-reduce & 定义它；写出一个最小例子；说明一个 failure mode；指出它影响哪个资源或评估账本。\\
条目 5 & communication cost（通信成本） & 定义它；写出一个最小例子；说明一个 failure mode；指出它影响哪个资源或评估账本。\\
\bottomrule
\end{longtable}
\subsection{综合自测}
\begin{itemize}
\item 我能否不用课程原句，独立解释本章中心问题？
\item 我能否把本章至少一个公式改写成可运行的估算函数？
\item 我能否指出一个看似改进但实际转移成本的方法？
\item 我能否回到视频和课程资料，找到支持本章关键论断的位置？
\item 我能否设计一个最小 ablation，验证本章一个 caveat？
\end{itemize}
\section{总结与延伸}
\subsection{本章小结}
\begin{itemize}
\item 多 GPU 并行基础 的核心不是记忆术语，而是理解对象、变量和机制之间的关系。
\item 视频给出讲解顺序，课程资料提供可核对的公式、图、代码或页面结构。
\item 复习时应把每个术语连接到至少一个变量、一个实现动作和一个失败模式。
\end{itemize}
\subsection{延伸解释}
本章中的延伸解释用于连接课程材料中的概念，例如说明公式里的 proxy、解释系统 tradeoff，或把多个材料片段串成一个完整推理链。若需要严肃引用，应回到课程视频和配套资料核对原始表述。
\subsection{练习}
\begin{itemize}
\item 定义题：用中英双语解释本章任意五个重要术语，并说明它们属于 compute、memory、data、optimization、evaluation 还是 deployment 账本。
\item 推导题：选择本章一个公式，逐个解释符号、单位和近似假设，并说明哪个变量最难在真实系统中测量。
\item 实现题：把本章伪代码改写成一个最小 Python sanity check，不要求训练大模型，但要输出可检查的中间量。
\item 资料题：从“课程资料与回看路径”中选择一个视频时间段，回到原始视频核对讲者例子，并记录是否存在字幕误差。
\item 评估题：为本章方法设计一个 ablation 或 benchmark，说明控制变量、数据版本、指标和 failure criteria。
\item 边界题：说明本章一个结论在更长 context、更大 batch、更多 GPU 或更复杂数据分布下可能如何失效。
\item 综合题：把本章内容和前一章或后一章连接起来，写出一个端到端训练/推理 pipeline 中的依赖关系。
\item 批判题：指出一个容易被营销材料夸大的说法，并用本章的资源账本或评估协议反驳它。
\end{itemize}
\end{document}
